\documentclass[12pt,a4paper]{article}

\usepackage[utf8]{inputenc}
\usepackage[T1]{fontenc}
\usepackage[brazilian]{babel}
\usepackage[margin=2.5cm]{geometry}
\usepackage{times}
\usepackage{parskip}

\begin{document}

{\LARGE\bfseries Apêndice A --- Declaração de uso de ferramentas de Inteligência Artificial}

\vspace{0.5cm}

Eu, Felipe Echeverria Vilhalva, declaro que utilizei ferramentas de
Inteligência Artificial (IA) no processo de elaboração do trabalho intitulado
``Arquitetura de Análise de Redes Parlamentares Baseada em Teoria dos
Grafos para Identificação de Estruturas de Influência'', conforme detalho
abaixo:

\vspace{0.3cm}
1. Ferramentas utilizadas:

NotebookLM (Google) e Claude Code (Anthropic).

\vspace{0.3cm}
2. Etapas em que foram utilizadas:

NotebookLM: apoio ao levantamento bibliográfico, auxiliando na identificação
e organização preliminar de fontes e referências relacionadas aos temas do
trabalho. A busca em bases reais e a verificação de cada fonte citada foram
realizadas pelo autor.

Claude Code: apoio à revisão gramatical e à reescrita de trechos para
aprimorar coesão e coerência, sempre a partir de texto previamente redigido
pelo autor; discussão e esclarecimento de conceitos técnicos; geração,
otimização e refatoração de código-fonte, incluindo os scripts de geração
dos gráficos e figuras analíticas (bibliotecas matplotlib/seaborn)
apresentados no Capítulo de Resultados.

\vspace{0.3cm}
3. Em quais seções do manuscrito:

Revisão gramatical e reescrita para coesão e coerência ao longo de todos os
capítulos, do Capítulo 1 ao Capítulo 6. Apoio bibliográfico concentrado na
Revisão da Literatura e na Metodologia. Apoio técnico (código-fonte, scripts
de análise e geração de gráficos) concentrado no Capítulo de Desenvolvimento
e nos utilitários de reprodutibilidade (scripts/).

\vspace{0.3cm}
4. Descrição dos limites de uso:

O uso foi restrito a apoio redacional (clareza, gramática e reescrita de
trechos já elaborados pelo autor), discussão de conceitos técnicos e
assistência na geração e refatoração de código. A concepção da arquitetura de
software, as decisões metodológicas e de modelagem, a validação lógica do
sistema, a interpretação dos resultados e a autoria intelectual do trabalho
são de responsabilidade exclusiva do autor, que verificou integralmente
todas as referências, dados e informações apresentados. Todo código
produzido com apoio das ferramentas foi submetido à suíte de testes
automatizados do artefato, e todos os resultados gráficos foram validados de
forma independente pelo autor.

\vspace{0.5cm}

Declaro ainda que assumo total responsabilidade pelo conteúdo do manuscrito,
garantindo sua originalidade, precisão, validade e integridade ética.

\vspace{0.8cm}

Local e Data: Dourados, MS, 15 de outubro de 2026.

\vspace{2cm}

\noindent\rule{7cm}{0.4pt}

\noindent Assinatura do(a) autor(a)

\vspace{1.5cm}

\noindent\rule{7cm}{0.4pt}

\noindent Assinatura do(a) orientador(a)

\end{document}
